%% Original pipeline figures removed from pediatric.tex on 2026-09-13 at the
%% author's request while real versions are drawn. Paste either block back in
%% place of its placeholder to restore.

% ---- original fig:architecture ----
\begin{figure*}[t]
\centering
\begin{tikzpicture}[
  node distance=7mm and 7mm,
  block/.style={draw,rounded corners,minimum height=9mm,align=center,fill=lightgray,font=\small},
  general/.style={block,fill=ctblue!13,draw=ctblue},
  context/.style={block,fill=ctxorange!15,draw=ctxorange},
  output/.style={block,fill=ctxgreen!14,draw=ctxgreen},
  arrow/.style={-{Latex[length=2mm]},thick}
]
\node[general] (ct) {3-D CT\\volume};
\node[general,right=of ct] (enc) {R(2+1)D-18 encoder\\visual tokens};
\node[general,right=of enc] (qgen) {General Q-Former\\10 anatomy + 23 pathology\\+ 2 global};
\node[context,below=of enc] (meta) {Indication + age band + sex\\clinical context $H_C$};
\node[context,right=of qgen] (c1) {C1 conditioned bank\\23 pathology twins\\+ 1 clinical};
\node[output,right=of c1,yshift=7mm] (prompt) {Conditioned\\prompt logits};
\node[output,right=of c1,yshift=-7mm] (cls) {Supervised\\\texttt{ctx\_cls} logits};
\draw[arrow] (ct)--(enc);
\draw[arrow] (enc)--(qgen);
\draw[arrow] (meta)-|(c1);
\draw[arrow] (qgen)--(c1);
\draw[arrow] (c1)--(prompt);
\draw[arrow] (qgen.east) to[out=15,in=165] (prompt.west);
\draw[arrow] (c1)--(cls);
\draw[arrow] (qgen.east) to[out=-15,in=165] (cls.west);
\node[draw,dashed,rounded corners,fit=(qgen)(c1),inner sep=3mm,label=above:{metadata-conditioned Q-Former}] {};
\end{tikzpicture}
\caption{Architecture for the 23-finding pediatric task. The
metadata-independent general bank is retained; pathology-specific conditioned
twins attend to clinical context through C1 cross-attention and FiLM and are
pooled by C2 relevance weights. Prompt inference is the primary readout;
\texttt{ctx\_cls} is a separate supervised ablation.}
\label{fig:architecture}
\end{figure*}

% ---- original fig:cohort-flow ----
\begin{figure*}[t]
\centering
\begin{tikzpicture}[
  node distance=6mm and 19mm,
  flow/.style={draw,rounded corners,minimum height=9mm,minimum width=28mm,
    align=center,fill=lightgray,font=\small},
  keep/.style={flow,fill=ctblue!12,draw=ctblue},
  split/.style={flow,fill=ctxgreen!12,draw=ctxgreen},
  arrow/.style={-{Latex[length=2mm]},thick}
]
\node[flow] (raw) {Institutional archive\\8,870 records};
\node[keep,right=of raw] (clean) {Model-ready CTs\\8,760 studies};
\node[keep,right=of clean] (u18) {$0\leq$age$<18$\\6,341 studies};
\node[split,right=16mm of u18,yshift=15mm] (train) {Train\\4,320 studies\\2,108 patients};
\node[split,right=16mm of u18] (dev) {Development\\748 studies\\373 patients};
\node[split,right=16mm of u18,yshift=-15mm] (test) {Frozen test\\1,273 studies\\611 patients};
\draw[arrow] (raw)--node[midway,above=1pt,font=\scriptsize]{conversion/QC}(clean);
\draw[arrow] (clean)--node[midway,above=1pt,font=\scriptsize]{age eligibility}(u18);
\draw[arrow] (u18)--(train.west);
\draw[arrow] (u18)--(dev.west);
\draw[arrow] (u18)--(test.west);
\end{tikzpicture}
\caption{Assembly of the primary pediatric cohort. Train, development, and
test are mutually exclusive at the MRN/patient level. The historical source
validation pool is retained as the frozen test cohort; the source training
pool alone is divided into train and development.}
\label{fig:cohort-flow}
\end{figure*}
